\section{Method}
\label{sec:method}

\subsection{Three-stage training design}
Our pipeline has three sequential stages, summarised in
\cref{tab:stages}.

\begin{table}[h]
\centering
\small
\begin{tabularx}{\linewidth}{lXrl}
\toprule
\textbf{Stage} & \textbf{Purpose} & \textbf{Train size} & \textbf{Data} \\
\midrule
S1: base reference     & Provide a starting score with no training            & 0       & --- \\
S2: SFT               & Teach answer format, reasoning style, output schema  & 1\,298  & \code{strict\_main} \\
S3: GRPO/DAPO         & Optimise verifiable correctness on a large pool      & up to 81\,463 & \code{rl\_strict} or \code{sampled24k} \\
\bottomrule
\end{tabularx}
\caption{Three training stages used in this report. S1 is evaluation only.}
\label{tab:stages}
\end{table}

\subsection{Base model}
We use \modelname{Qwen2.5-1.5B}~\cite{qwen25} as the base. Earlier internal
tests on a smaller \modelname{Qwen2.5-0.5B} variant proved unstable across
both SFT and RL, so all numbers in this report use the 1.5B variant. We
keep three model artefact types throughout the project, listed in
\cref{tab:models}.

\begin{table}[h]
\centering
\small
\begin{tabularx}{\linewidth}{lXX}
\toprule
\textbf{Form} & \textbf{Identifier} & \textbf{Use} \\
\midrule
\code{base}        & \modelname{Qwen2.5-1.5B}                                 & Reference for S1 \\
\code{sft}         & \code{outputs/sft/strict\_main\_qwen25\_15b\_base\_bs32\_ga1}    & LoRA SFT result, S2 output \\
\code{sft merged}  & \code{qwen base sft merged}                              & Full-parameter base for GRPO/verl \\
\bottomrule
\end{tabularx}
\caption{Three model forms used during training and evaluation.}
\label{tab:models}
\end{table}

\subsection{Stage 2: SFT}
The S2 stage performs supervised fine-tuning on the
\code{strict\_main} ShareGPT export. We use a LoRA adapter with rank 16,
batch size 32, and one gradient-accumulation step on a single GPU. The
training objective is the standard next-token cross-entropy on the
assistant turn only. The aim is purely to teach the model:

\begin{itemize}
  \item the math instruction template (problem $+$ requirements);
  \item the structure of a concise step-by-step reasoning trace;
  \item the unique \code{Final Answer: <answer>} closing line.
\end{itemize}

After S2, the LoRA adapter can either be evaluated directly (rows ``qwen
base sft'' below) or merged into the base weights to provide a starting
point for full-parameter RL.

\subsection{Stage 3: GRPO and DAPO}
Stage 3 applies a verifiable-reward RL algorithm. We use two related
algorithms:

\begin{itemize}
  \item \textbf{GRPO}~\cite{deepseekmath}, which estimates the advantage by
    standardising rewards within a group of $G$ rollouts that share the
    same prompt, removing the need for a separate value network.
  \item \textbf{DAPO}~\cite{dapo}, which extends GRPO with decoupled
    clipping bounds and a dynamic-sampling strategy that drops trivial or
    impossible prompts on the fly. We use the open-source \emph{verl}
    framework~\cite{verl} to run both algorithms on a single GPU.
\end{itemize}

For a prompt $q$ with reference answer $a^{*}$ and a group of $G$ sampled
completions $\{o_i\}_{i=1}^{G}$, GRPO uses the standardised advantage
\begin{equation}
\hat A_i \;=\; \frac{R(q, o_i, a^{*}) - \mu_R}{\sigma_R + \varepsilon},
\qquad i = 1,\dots,G,
\label{eq:advantage}
\end{equation}
where $\mu_R$ and $\sigma_R$ are the mean and standard deviation of the
group rewards, $\varepsilon$ is a small constant for numerical stability,
and $R(\cdot)$ is the composite reward defined in \cref{sec:reward}.

The policy is then updated with a PPO-style clipped objective
\begin{equation}
\mathcal{L}_{\text{policy}}(\theta) \;=\;
- \mathbb{E}\!\left[\min\!\Big(\rho_i\, \hat A_i,\;
  \mathrm{clip}(\rho_i, 1-\epsilon,\, 1+\epsilon)\, \hat A_i\Big)\right],
\quad
\rho_i \;=\; \frac{\pi_\theta(o_i \mid q)}{\pi_{\theta_{\text{old}}}(o_i \mid q)}.
\label{eq:ppo}
\end{equation}
A KL term against the SFT reference policy is added with a small
coefficient. The DAPO variant replaces the symmetric clip $\epsilon$ with
the asymmetric pair $(\epsilon_{\text{low}}, \epsilon_{\text{high}})$ and
filters prompts whose group reward is degenerate.

\subsection{Checkpoints reported}
Within stage 3 we evaluate three named checkpoints:
\code{grpo sampled24k}, \code{dapo}, and \code{grpo}. The first uses the
24K balanced sub-sample, the latter two use the full
\code{rl\_strict} pool of 81\,463 problems. Both \code{dapo} and \code{grpo}
correspond to mid-training checkpoints (step 150) rather than the very last
training step, since their accuracy curves had already plateaued.
